% Generated by scripts/paper/fig_method.R; do not edit by hand.
\begin{table}[tbp]
  \centering
  \small
  \begin{threeparttable}
    \caption{Components of the Policy index}\label{tab:method_policy}
    \begin{tabular}{@{}llrl@{}}
      \toprule
      Component & Weight & Economies & In the composite \\
      \midrule
      IEA Policies and Measures (PAMS) & 3 & 215 & yes, required \\
      Climate Action Tracker rating & 2 & 65 & where rated (39 own, 26 via the EU) \\
      GTA New Industrial Policy Observatory & -- & 77 & no (robustness variant, October 2026) \\
      Dual-use policy score & -- & 217 & no (constant within pairs) \\
      \midrule
      Economies with a Policy index & & 215 & \\
      \quad on IEA PAMS and CAT & & 65 & \\
      \quad on IEA PAMS only & & 150 & CAT weight renormalised away \\
      Coalition members with a Policy index & & 57 & of 57 (11 on IEA PAMS only) \\
      \bottomrule
    \end{tabular}
    \begin{tablenotes}[flushleft]\footnotesize
      \item \emph{Notes:} Economies: economies with at least one value of the component in the
      eleven headline technologies. The Policy index of a cell is the weighted mean of the
      composite components present in it, with the weights renormalised over those components.
      IEA PAMS is required: an economy is scored when it has an IEA PAMS value in at least one
      headline cell, and the Climate Action Tracker rating is added where the economy has one.
      EU member states without a rating of their own take the EU's rating. Requiring both
      components (\texttt{OPSI\_POLICY\_REQUIRE\_ALL=1}, a robustness variant) leaves the 65 economies scored on both. For those 65, the within-country Spearman correlation between the IEA PAMS component alone and the composite over their headline cells has median 0.48 (interquartile range 0.41--0.57), so the basis on which an economy is scored affects the ordering of its sectors.
      The NIPO component enters only a robustness variant, built from the October 2026 NIPO release (interventions in force on 5 October 2026; 77 jurisdictions, EU members individually); the dual-use score takes one value for every economy of a technology--stage pair, so it cannot rank economies.
      Source: \texttt{policy\_component\_tbl}, \texttt{policy\_eligibility}, \texttt{policy\_index} and
      \texttt{coalition\_members} in \texttt{data/processed/outputs/index\_outputs.rds};
      \texttt{scripts/paper/fig\_method.R}.
    \end{tablenotes}
  \end{threeparttable}
\end{table}
